\begin{abstract}
异构图神经网络能够建模包含多类型节点和多类型边的复杂关系数据，已经在节点分类和表示学习任务中取得了较好效果。
然而，这类模型仍然容易受到结构对抗攻击影响：少量恶意边扰动可能通过基于元路径的消息传播被放大，并进一步污染高阶语义邻域。
现有鲁棒异构图方法通常在拓扑语义空间内进行防御，例如估计边置信度、过滤可疑的元路径诱导边，或学习语义注意力权重。
这些方法在一定程度上有效，但仍然高度依赖被攻击后的异构拓扑来构造候选邻域；当拓扑本身已被严重污染时，模型很难恢复可靠表示。

为了解决这一问题，本文提出 DVCL，一种面向鲁棒异构图学习的双视图对比学习框架。
DVCL 构造两个互补视图：由元路径诱导语义结构得到的净化拓扑视图，以及由节点特征相似性构建的特征诱导视图。
拓扑视图通过元路径建模、语义注意力和两阶段边过滤保留异构结构语义；特征视图则在结构攻击下提供相对独立于拓扑的参考。
为了降低两个视图之间的语义不一致，DVCL 使用对称跨视图对比目标，使同一节点的拓扑表示与特征表示相互对齐，同时区分不同节点的表示。
结合监督节点分类目标和辅助语义拓扑学习目标，DVCL 学习兼具异构语义表达能力与特征层稳定性的鲁棒节点表示。
在结构攻击下的异构图基准实验表明，相比现有异构图防御基线，DVCL 能够提升节点分类鲁棒性。
\end{abstract}
